\documentclass{article}
\usepackage[utf8]{inputenc}
\usepackage[T1]{fontenc}
\usepackage{amsmath}
\usepackage{amssymb}
\usepackage{mathtools}
\usepackage[margin=1in]{geometry}
\usepackage{enumitem}
\usepackage{parskip}
\usepackage{array}
\usepackage{booktabs}
\usepackage{hyperref}

\title{CEC 315 --- Exam 3 Study Guide (Lectures 16--23)}
\author{Course: CEC 315 --- Signals and Systems (Spring 2026)\\
Instructor: Rogelio Gracia Otalvaro\\
Coverage: Lectures 16--23 (Laplace, Z-Transform, Sampling, Feedback)}
\date{}

\begin{document}
\maketitle

\begin{center}
\fbox{\parbox{0.92\textwidth}{%
\textbf{\large AUTHORITATIVE SOURCE: Official Instructor Study Guide}\\[4pt]
The instructor released an \textbf{official Exam 3 study guide} (\texttt{StudyGuide\_Exam3.pdf}). A faithful transcription is in \texttt{official\_study\_guide.md} (LaTeX: \texttt{official\_study\_guide.tex}). \textbf{Treat that file as authoritative} --- this companion guide is supplementary.\\[4pt]
\textbf{Exam format (from official guide):} 50 minutes, 100 pts total. Part I: 10 multiple choice ($4\times 10 = 40$ pts). Part II: multi-part problems (60 pts).\\[4pt]
\textbf{Instructor's \#1 emphasis:} \emph{``Always state the ROC with every transform.''}\\[4pt]
The official guide also includes \textbf{25 practice questions} (15 sampling, 10 feedback) --- review them before the exam.
}}
\end{center}

\tableofcontents
\newpage

\section{Topic Overview by Lecture}

\subsection*{Lecture 16 --- Laplace Transform and ROC (\S9.0--9.2)}

\noindent\emph{Instructor emphasis (official guide):} ROC is a vertical strip (depends only on $\sigma$), never contains poles, and $\dfrac{1}{s+a}$ is \emph{ambiguous without the ROC} --- could be $e^{-at}u(t)$ or $-e^{-at}u(-t)$.

\begin{itemize}
  \item Motivation: Fourier transform fails for growing signals (e.g., $e^{2t}u(t)$).
  \item Definition of the bilateral Laplace transform $X(s) = \int_{-\infty}^{\infty} x(t) e^{-st}\,dt$.
  \item The complex variable $s = \sigma + j\omega$ and the $s$-plane.
  \item Region of convergence (ROC): set of $s$ where the integral converges.
  \item Basic pairs: $e^{-at}u(t)$, $-e^{-at}u(-t)$, $\delta(t)$, $u(t)$, $t e^{-at}u(t)$.
  \item Poles, zeros, and pole-zero plots; ROC bounded by poles.
  \item Properties of the ROC (strip parallel to $j\omega$-axis, no poles, finite duration, right/left/two-sided rules).
  \item Fourier transform as a special case: $X(j\omega) = X(s)\big|_{s=j\omega}$ when the $j\omega$-axis lies in the ROC.
\end{itemize}

\subsection*{Lecture 17 --- Inverse Laplace and Properties (\S9.3--9.6)}
\begin{itemize}
  \item Partial fraction expansion (distinct, repeated, and complex conjugate poles).
  \item Using ROC to assign right-sided vs.\ left-sided terms.
  \item Geometric evaluation of $|H(j\omega)|$ from the pole-zero plot.
  \item Laplace properties: linearity, time shift, $s$-domain shift, time scaling, conjugation, convolution, differentiation in $t$, integration in $t$, differentiation in $s$.
  \item Initial-value theorem: $x(0^+) = \lim_{s\to\infty} s X(s)$.
  \item Final-value theorem: $x(\infty) = \lim_{s\to 0} s X(s)$ (only when valid).
\end{itemize}

\subsection*{Lecture 18 --- System Analysis via Unilateral Laplace (\S9.7--9.9)}

\noindent\emph{Instructor ``Golden Rule'' (CT) from official guide:} Causal + BIBO stable $\iff$ all poles in the open LHP ($\operatorname{Re}\{p_i\} < 0\ \forall i$). Unilateral differentiation uses \textbf{$-y(0^-)$}.

\begin{itemize}
  \item Deriving $H(s) = Y(s)/X(s)$ from a differential equation.
  \item Causality $\Leftrightarrow$ ROC is a right half-plane (and $M \le N$).
  \item BIBO stability $\Leftrightarrow$ ROC includes the $j\omega$-axis.
  \item ``Golden Rule'': causal + stable $\Leftrightarrow$ all poles strictly in the open LHP.
  \item Block diagrams: series ($H_1 H_2$), parallel ($H_1 + H_2$), feedback ($G/(1+GF)$).
  \item Unilateral Laplace: $\mathcal{X}(s) = \int_{0^-}^\infty x(t)e^{-st}dt$.
  \item Differentiation rule with ICs: $\mathcal{L}_u\{x'\} = sX(s) - x(0^-)$.
  \item Solving differential equations with nonzero initial conditions.
  \item Zero-state response (ZSR) + zero-input response (ZIR) decomposition.
\end{itemize}

\subsection*{Lecture 19 --- Z-Transform and ROC (\S10.0--10.2)}
\begin{itemize}
  \item Motivation: DTFT fails for growing sequences (e.g., $2^n u[n]$).
  \item Definition: $X(z) = \sum_n x[n] z^{-n}$; $z = re^{j\omega}$.
  \item Unit circle $|z|=1$ plays the role of the $j\omega$-axis.
  \item Basic pairs: $a^n u[n]$, $-a^n u[-n-1]$, $\delta[n]$, $\delta[n-k]$, $u[n]$.
  \item Right-sided $\Rightarrow$ ROC outside outermost pole; left-sided $\Rightarrow$ ROC inside innermost pole; two-sided $\Rightarrow$ annular ring.
  \item Finite-duration $\Rightarrow$ entire $z$-plane (excluding possibly $z=0$ or $\infty$).
  \item DTFT exists $\Leftrightarrow$ unit circle in ROC.
\end{itemize}

\subsection*{Lecture 20 --- Inverse Z-Transform and Properties (\S10.3--10.6)}
\begin{itemize}
  \item Partial fractions in $z^{-1}$ (crucial: not in $z$!).
  \item Right-sided vs.\ left-sided direction assignment from the ROC.
  \item Complex poles at $re^{\pm j\omega_0}$ give damped sinusoids $r^n\cos(\omega_0 n)$.
  \item Geometric evaluation of frequency response on the unit circle.
  \item Z properties: linearity, time shift, z-domain scaling, time reversal, convolution, differentiation in $z$, first difference, accumulation.
  \item Initial-value theorem: $x[0] = \lim_{z\to\infty} X(z)$.
  \item Final-value theorem: $x[\infty] = \lim_{z\to 1}(1-z^{-1})X(z)$.
\end{itemize}

\subsection*{Lecture 21 --- System Analysis via Unilateral Z (\S10.7--10.9)}

\noindent\emph{Instructor ``Sign trap'' (official guide):} Laplace uses $-y(0^-)$; the $z$-transform uses \textbf{$+y[-1]$}. Don't mix them up. \emph{Golden Rule (DT):} causal + BIBO stable $\iff$ all poles strictly inside the unit circle ($|p_i| < 1$).

\begin{itemize}
  \item Deriving $H(z) = Y(z)/X(z)$ from a difference equation (replace $y[n-k] \to z^{-k}Y(z)$).
  \item Causality $\Leftrightarrow$ ROC is exterior of a circle; $M \le N$.
  \item BIBO stability $\Leftrightarrow$ ROC includes the unit circle.
  \item ``Golden Rule'' (DT): causal + stable $\Leftrightarrow$ all poles strictly inside the unit circle ($|p_i|<1$).
  \item Block diagrams with $z^{-1}$ as the unit delay.
  \item Unilateral z-transform: $\mathcal{X}(z) = \sum_{n=0}^\infty x[n]z^{-n}$.
  \item Shift rule with ICs: $x[n-1] \leftrightarrow z^{-1}X(z) + x[-1]$ (note the ``+'', not ``$-$'' like Laplace).
  \item Solving difference equations with initial conditions; ZSR + ZIR.
\end{itemize}

\subsection*{Lecture 22 --- Sampling (Ch.\ 7)}

\noindent\emph{Instructor emphasis:} ``Sampling = copy-paste the spectrum at every multiple of $\omega_s$, scaled by $1/T$.'' $\omega_s = 2\omega_M$ is \textbf{not sufficient} (strict $>$ required; e.g.\ $\sin(\omega_s t/2)$ sampled at $\omega_s$ gives all-zero samples). \textbf{Squaring a signal doubles its bandwidth.} The official guide includes 15 practice questions for this lecture --- work through them.

\begin{itemize}
  \item Impulse-train sampling: $x_p(t) = x(t) p(t)$, $p(t) = \sum_n \delta(t-nT)$.
  \item Frequency-domain effect: spectrum replicated at multiples of $\omega_s$, scaled by $1/T$.
  \item Sampling theorem: $\omega_s > 2\omega_M$ for perfect reconstruction (strict inequality!).
  \item Ideal reconstruction: lowpass filter with gain $T$, cutoff $\omega_c$ between $\omega_M$ and $\omega_s-\omega_M$.
  \item Band-limited interpolation (sinc interpolation formula).
  \item Practical interpolation: zero-order hold (staircase), first-order hold (linear).
  \item Aliasing: $f_\text{alias} = |f_s - f_0|$ when undersampled.
  \item Wagon-wheel effect, moir\'e patterns.
  \item Discrete-time processing of CT signals: $\Omega = \omega T$.
\end{itemize}

\subsection*{Lecture 23 --- Linear Feedback Systems (\S11.0--11.5)}

\noindent\emph{Instructor emphasis (official guide):} Know the three regimes of $|GH|$ --- $|GH|\gg 1 \Rightarrow Q\approx 1/G$ (feedback dominates); $|GH|\ll 1 \Rightarrow Q\approx H$ (feedback negligible); $GH=-1 \Rightarrow$ instability. \emph{Strategy:} simplify innermost loop first. \emph{Reading Bode plots:} GM from the \emph{magnitude} plot at the $-180^\circ$ frequency; PM from the \emph{phase} plot at the 0 dB frequency --- don't mix them up! Official guide has 10 feedback practice questions (block diagram simplification, closed-loop stability, GM/PM).

\begin{itemize}
  \item Open-loop vs.\ closed-loop system structure.
  \item Closed-loop transfer function: $Q(s) = H/(1+GH)$ for negative feedback.
  \item Block diagram simplification: cascade, parallel, feedback.
  \item Closed-loop poles from the characteristic equation $1 + GH = 0$.
  \item Effect of feedback: disturbance rejection, sensitivity reduction, pole relocation.
  \item Nyquist plot: polar plot of $G(j\omega)H(j\omega)$.
  \item Nyquist stability criterion (simplified): closed-loop stable iff Nyquist plot does not encircle $-1/K$ (given open-loop is stable).
  \item Gain margin and phase margin from Bode plots.
  \item CT vs.\ DT feedback stability boundaries.
\end{itemize}

\section{Key Equations Reference}

\subsection{Transform Definitions}

\begin{center}
\begin{tabular}{ll}
\toprule
Transform & Formula \\
\midrule
Bilateral Laplace & $\displaystyle X(s) = \int_{-\infty}^{\infty} x(t)e^{-st}\,dt$ \\
Unilateral Laplace & $\displaystyle \mathcal{X}(s) = \int_{0^-}^{\infty} x(t)e^{-st}\,dt$ \\
Inverse Laplace & $\displaystyle x(t) = \frac{1}{2\pi j}\int_{\sigma-j\infty}^{\sigma+j\infty} X(s)e^{st}\,ds$ \\
Bilateral Z & $\displaystyle X(z) = \sum_{n=-\infty}^{\infty} x[n]z^{-n}$ \\
Unilateral Z & $\displaystyle \mathcal{X}(z) = \sum_{n=0}^{\infty} x[n]z^{-n}$ \\
Inverse Z & $\displaystyle x[n] = \frac{1}{2\pi j}\oint X(z)z^{n-1}\,dz$ \\
\bottomrule
\end{tabular}
\end{center}

\subsection{Common Laplace Pairs (all causal, ROC = right of rightmost pole)}

\begin{center}
\begin{tabular}{lll}
\toprule
$x(t)$ & $X(s)$ & ROC \\
\midrule
$\delta(t)$ & $1$ & all $s$ \\
$\delta(t-t_0)$ & $e^{-st_0}$ & all $s$ \\
$u(t)$ & $\dfrac{1}{s}$ & $\operatorname{Re}\{s\}>0$ \\
$tu(t)$ & $\dfrac{1}{s^2}$ & $\operatorname{Re}\{s\}>0$ \\
$\dfrac{t^{n-1}}{(n-1)!}u(t)$ & $\dfrac{1}{s^n}$ & $\operatorname{Re}\{s\}>0$ \\
$e^{-at}u(t)$ & $\dfrac{1}{s+a}$ & $\operatorname{Re}\{s\}>-\operatorname{Re}\{a\}$ \\
$t e^{-at}u(t)$ & $\dfrac{1}{(s+a)^2}$ & $\operatorname{Re}\{s\}>-\operatorname{Re}\{a\}$ \\
$\dfrac{t^{n-1}}{(n-1)!}e^{-at}u(t)$ & $\dfrac{1}{(s+a)^n}$ & $\operatorname{Re}\{s\}>-\operatorname{Re}\{a\}$ \\
$\cos(\omega_0 t)u(t)$ & $\dfrac{s}{s^2+\omega_0^2}$ & $\operatorname{Re}\{s\}>0$ \\
$\sin(\omega_0 t)u(t)$ & $\dfrac{\omega_0}{s^2+\omega_0^2}$ & $\operatorname{Re}\{s\}>0$ \\
$e^{-at}\cos(\omega_d t)u(t)$ & $\dfrac{s+a}{(s+a)^2+\omega_d^2}$ & $\operatorname{Re}\{s\}>-a$ \\
$e^{-at}\sin(\omega_d t)u(t)$ & $\dfrac{\omega_d}{(s+a)^2+\omega_d^2}$ & $\operatorname{Re}\{s\}>-a$ \\
$-e^{-at}u(-t)$ & $\dfrac{1}{s+a}$ & $\operatorname{Re}\{s\}<-\operatorname{Re}\{a\}$ \\
\bottomrule
\end{tabular}
\end{center}

\subsection{Common Z-Transform Pairs}

\begin{center}
\begin{tabular}{lll}
\toprule
$x[n]$ & $X(z)$ & ROC \\
\midrule
$\delta[n]$ & $1$ & all $z$ \\
$\delta[n-k]$, $k\ge 0$ & $z^{-k}$ & $|z|>0$ \\
$u[n]$ & $\dfrac{1}{1-z^{-1}}$ & $|z|>1$ \\
$a^n u[n]$ & $\dfrac{1}{1-az^{-1}}$ & $|z|>|a|$ \\
$-a^n u[-n-1]$ & $\dfrac{1}{1-az^{-1}}$ & $|z|<|a|$ \\
$na^n u[n]$ & $\dfrac{az^{-1}}{(1-az^{-1})^2}$ & $|z|>|a|$ \\
$(n+1)a^n u[n]$ & $\dfrac{1}{(1-az^{-1})^2}$ & $|z|>|a|$ \\
$r^n \cos(\omega_0 n)u[n]$ & $\dfrac{1-r\cos\omega_0\, z^{-1}}{1-2r\cos\omega_0\, z^{-1}+r^2 z^{-2}}$ & $|z|>r$ \\
$r^n \sin(\omega_0 n)u[n]$ & $\dfrac{r\sin\omega_0\, z^{-1}}{1-2r\cos\omega_0\, z^{-1}+r^2 z^{-2}}$ & $|z|>r$ \\
\bottomrule
\end{tabular}
\end{center}

\subsection{Laplace Properties}

\begin{center}
\begin{tabular}{lll}
\toprule
Property & Time & $s$-Domain \\
\midrule
Linearity & $ax(t)+by(t)$ & $aX(s)+bY(s)$, ROC $\supseteq R_x\cap R_y$ \\
Time shift & $x(t-t_0)$ & $e^{-st_0}X(s)$, ROC $=R_x$ \\
$s$-domain shift & $e^{s_0 t}x(t)$ & $X(s-s_0)$, ROC $=R_x+\operatorname{Re}\{s_0\}$ \\
Time scaling & $x(at)$, $a>0$ & $\dfrac{1}{a}X(s/a)$, ROC $=aR_x$ \\
Conjugation & $x^*(t)$ & $X^*(s^*)$, ROC $=R_x$ \\
Convolution & $x(t)*y(t)$ & $X(s)Y(s)$ \\
Differentiation in $t$ & $dx/dt$ & $sX(s)$ (bilateral, zero ICs) \\
Integration in $t$ & $\int_{-\infty}^t x(\tau)d\tau$ & $X(s)/s$ \\
Differentiation in $s$ & $-tx(t)$ & $dX/ds$ \\
\textbf{Unilateral diff.} & $dx/dt$ & $sX(s)-x(0^-)$ \\
Unilateral $d^2/dt^2$ & $d^2x/dt^2$ & $s^2X(s)-sx(0^-)-x'(0^-)$ \\
Initial-value theorem & $x(0^+)$ & $\lim_{s\to\infty} sX(s)$ \\
Final-value theorem & $x(\infty)$ & $\lim_{s\to 0} sX(s)$ (if valid) \\
\bottomrule
\end{tabular}
\end{center}

\subsection{Z-Transform Properties}

\begin{center}
\begin{tabular}{lll}
\toprule
Property & Sequence & z-Domain \\
\midrule
Linearity & $ax_1+bx_2$ & $aX_1+bX_2$ \\
Time shift & $x[n-n_0]$ & $z^{-n_0}X(z)$ \\
z-domain scaling & $z_0^n x[n]$ & $X(z/z_0)$, ROC $=|z_0|R_x$ \\
Time reversal & $x[-n]$ & $X(z^{-1})$, ROC $=1/R_x$ \\
Convolution & $x_1*x_2$ & $X_1(z)X_2(z)$ \\
Differentiation in $z$ & $nx[n]$ & $-z\,dX/dz$ \\
First difference & $x[n]-x[n-1]$ & $(1-z^{-1})X(z)$ \\
Accumulation & $\sum_{k=-\infty}^n x[k]$ & $\dfrac{X(z)}{1-z^{-1}}$ \\
\textbf{Unilateral shift} & $x[n-1]$ & $z^{-1}X(z)+x[-1]$ \\
Unilateral shift-2 & $x[n-2]$ & $z^{-2}X(z)+x[-1]z^{-1}+x[-2]$ \\
Initial-value theorem & $x[0]$ & $\lim_{z\to\infty} X(z)$ \\
Final-value theorem & $x[\infty]$ & $\lim_{z\to 1}(1-z^{-1})X(z)$ \\
\bottomrule
\end{tabular}
\end{center}

\section{ROC Rules (Laplace and Z)}

\subsection{Laplace ROC}

\begin{center}
\begin{tabular}{ll}
\toprule
Signal type & ROC shape \\
\midrule
Finite duration (absolutely integrable) & Entire $s$-plane \\
Right-sided ($x(t)=0$ for $t<T_0$) & $\operatorname{Re}\{s\} > \sigma_{\max}$ (right of rightmost pole) \\
Left-sided ($x(t)=0$ for $t>T_0$) & $\operatorname{Re}\{s\} < \sigma_{\min}$ (left of leftmost pole) \\
Two-sided & Vertical strip between consecutive poles \\
Causal and stable & Rightmost pole in LHP; $j\omega$-axis in ROC \\
\bottomrule
\end{tabular}
\end{center}

\textbf{ROC key facts:}
\begin{itemize}
  \item ROC is always a \textbf{vertical strip} parallel to the $j\omega$-axis.
  \item ROC never contains poles.
  \item DTFT/Fourier transform exists iff ROC includes the $j\omega$-axis.
\end{itemize}

\subsection{Z-Transform ROC}

\begin{center}
\begin{tabular}{ll}
\toprule
Signal type & ROC shape \\
\midrule
Finite duration & Entire $z$-plane (possibly excluding $z=0$ or $\infty$) \\
Right-sided (causal) & $|z| > r_{\max}$ (exterior of outermost pole) \\
Left-sided (anti-causal) & $|z| < r_{\min}$ (interior of innermost pole) \\
Two-sided & Annular ring $r_1 < |z| < r_2$ between pole circles \\
Causal and stable & All poles inside unit circle; unit circle in ROC \\
\bottomrule
\end{tabular}
\end{center}

\textbf{ROC key facts:}
\begin{itemize}
  \item ROC is always an \textbf{annular ring} centered at the origin.
  \item ROC never contains poles.
  \item DTFT exists iff ROC includes the unit circle.
\end{itemize}

\subsection{Causality and Stability Summary}

\begin{center}
\begin{tabular}{lll}
\toprule
Criterion & Laplace (CT) & Z (DT) \\
\midrule
Causal & ROC right half-plane, $M\le N$ & ROC exterior of circle, $M\le N$ \\
BIBO stable & ROC includes $j\omega$-axis & ROC includes unit circle \\
Causal + stable & All poles with $\operatorname{Re}\{p_i\}<0$ & All poles with $|p_i|<1$ \\
Marginally stable & Poles on $j\omega$-axis & Poles on unit circle \\
\bottomrule
\end{tabular}
\end{center}

\section{Sampling Theorem}

\subsection{The Theorem}

Let $x(t)$ be bandlimited with $X(j\omega)=0$ for $|\omega|>\omega_M$. Then $x(t)$ is uniquely determined by samples $x(nT)$ if and only if
$$\omega_s > 2\omega_M, \qquad \omega_s = \frac{2\pi}{T}.$$

\begin{itemize}
  \item \textbf{Nyquist rate:} $2\omega_M$ (minimum sampling frequency).
  \item \textbf{Nyquist frequency:} $\omega_M$.
  \item The inequality is \textbf{strict}: $\omega_s = 2\omega_M$ is generally not sufficient.
\end{itemize}

\subsection{Spectrum of Sampled Signal}

$$X_p(j\omega) = \frac{1}{T}\sum_{k=-\infty}^{\infty} X(j(\omega - k\omega_s))$$

The spectrum is replicated every $\omega_s$ and scaled by $1/T$.

\subsection{Ideal Reconstruction}

Use LPF with gain $T$ and cutoff $\omega_c$ such that $\omega_M < \omega_c < \omega_s - \omega_M$. Typical choice: $\omega_c = \omega_s/2$.

Time-domain formula:
$$x_r(t) = \sum_{n=-\infty}^{\infty} x(nT)\,\frac{\sin[\pi(t-nT)/T]}{\pi(t-nT)/T}$$

\subsection{Aliasing}

When $\omega_s < 2\omega_M$, the replicas overlap. A sinusoid at $f_0$ (with $f_s/2 < f_0 < f_s$) aliases to
$$f_\text{alias} = |f_s - f_0|.$$

\textbf{Aliasing is irreversible.} Prevent it with an \textbf{anti-aliasing filter} before sampling.

\subsection{DT Processing of CT Signals}

$$\Omega = \omega T$$

Relates discrete-time frequency $\Omega$ (rad/sample) to continuous-time frequency $\omega$ (rad/s).

\section{Feedback System Analysis}

\subsection{Closed-Loop Transfer Function}

For standard negative feedback:
$$Q(s) = \frac{Y(s)}{X(s)} = \frac{H(s)}{1 + G(s)H(s)}$$

\begin{itemize}
  \item \textbf{Forward path:} $H(s)$
  \item \textbf{Feedback path:} $G(s)$
  \item \textbf{Loop gain:} $L(s) = G(s)H(s)$
  \item \textbf{Characteristic equation:} $1 + G(s)H(s) = 0$ (closed-loop pole locations)
  \item Positive feedback variant: $Q = H/(1 - GH)$
\end{itemize}

\subsection{Block Diagram Rules}

\begin{center}
\begin{tabular}{ll}
\toprule
Configuration & Equivalent \\
\midrule
Cascade (series) & $H_1 \cdot H_2$ \\
Parallel & $H_1 + H_2$ \\
Negative feedback & $H/(1+GH)$ \\
\bottomrule
\end{tabular}
\end{center}

\subsection{Stability via Nyquist Criterion}

For an \textbf{open-loop stable} plant $GH$ (no RHP poles), the closed-loop system with gain $K$ is stable iff the Nyquist plot of $G(j\omega)H(j\omega)$ does \textbf{not encircle} the critical point $-1/K$.

\begin{itemize}
  \item Plot $GH$ for $\omega \in (-\infty, \infty)$.
  \item Locate $-1/K$ on the real axis.
  \item Critical point outside the closed contour $\Rightarrow$ stable.
\end{itemize}

\subsection{Gain and Phase Margins}

Instability occurs when $GH = -1$ (magnitude 1, phase $-180^\circ$ simultaneously).

\textbf{Gain margin (GM):} At $\omega_1$ where $\angle GH(\omega_1) = -180^\circ$,
$$\mathrm{GM}_{\mathrm{dB}} = -20\log_{10}|GH(\omega_1)|$$

\textbf{Phase margin (PM):} At $\omega_2$ where $|GH(\omega_2)| = 1$ (0 dB),
$$\mathrm{PM} = 180^\circ + \angle GH(\omega_2)$$

Both positive $\Rightarrow$ stable. Max additional delay: $\tau_{\max} = \mathrm{PM\ (rad)}/\omega_2$.

\subsection{CT vs.\ DT Stability}

\begin{center}
\begin{tabular}{lll}
\toprule
 & CT & DT \\
\midrule
Closed-loop & $H(s)/[1+G(s)H(s)]$ & $H(z)/[1+G(z)H(z)]$ \\
Stable region & $\operatorname{Re}\{s\}<0$ (open LHP) & $|z|<1$ (inside unit circle) \\
\bottomrule
\end{tabular}
\end{center}

\section{What to Memorize}

\textbf{Must-know transform pairs:}
$$\delta(t)\leftrightarrow 1,\qquad u(t)\leftrightarrow \tfrac{1}{s},\qquad e^{-at}u(t)\leftrightarrow \tfrac{1}{s+a}$$
$$te^{-at}u(t)\leftrightarrow \tfrac{1}{(s+a)^2},\qquad \cos\omega_0 t\,u(t)\leftrightarrow \tfrac{s}{s^2+\omega_0^2},\qquad \sin\omega_0 t\,u(t)\leftrightarrow \tfrac{\omega_0}{s^2+\omega_0^2}$$

\textbf{Must-know Z pairs:}
$$\delta[n]\leftrightarrow 1,\qquad u[n]\leftrightarrow \tfrac{1}{1-z^{-1}},\qquad a^n u[n]\leftrightarrow \tfrac{1}{1-az^{-1}}\ (|z|>|a|)$$
$$(n+1)a^n u[n]\leftrightarrow \tfrac{1}{(1-az^{-1})^2}$$

\textbf{Must-know rules:}
\begin{itemize}
  \item $dx/dt \leftrightarrow sX(s)-x(0^-)$ (unilateral).
  \item $x[n-1] \leftrightarrow z^{-1}X(z)+x[-1]$ (unilateral, note the \textbf{plus}).
  \item Convolution $\leftrightarrow$ multiplication.
  \item Causal + stable $\Leftrightarrow$ all poles in LHP (CT) or inside unit circle (DT).
  \item Sampling theorem: $\omega_s > 2\omega_M$.
  \item Closed-loop: $Q = H/(1+GH)$.
  \item Aliased frequency: $f_\text{alias} = |f_s - f_0|$.
\end{itemize}

\textbf{Must-know steps:}
\begin{enumerate}
  \item \textbf{Solving an ODE with ICs:} (a) unilateral Laplace, (b) plug in IC terms, (c) solve for $Y(s)$, (d) partial fractions, (e) invert.
  \item \textbf{Finding impulse response:} (a) compute $H(s)=Y/X$, (b) partial fractions, (c) use ROC to assign direction, (d) invert each term.
  \item \textbf{Stability check:} (a) find poles, (b) causal + LHP $\Rightarrow$ stable (CT); causal + inside unit circle $\Rightarrow$ stable (DT).
\end{enumerate}

\section{Common Pitfalls}

\begin{enumerate}
  \item \textbf{Omitting the ROC.} $X(s)=1/(s+3)$ could be $e^{-3t}u(t)$ or $-e^{-3t}u(-t)$. Always state the ROC.
  \item \textbf{Wrong direction in partial fractions.} Use the ROC, \emph{not} the sign of the pole, to choose right- vs.\ left-sided.
  \item \textbf{Incomplete partial fractions for repeated poles.} A pole of order $n$ requires $n$ terms: $\frac{B_1}{s-p},\ldots,\frac{B_n}{(s-p)^n}$.
  \item \textbf{Mishandling complex poles in Laplace.} Don't split conjugate pairs; complete the square and use cos/sin pairs.
  \item \textbf{Forgetting IC terms in unilateral differentiation.} $\mathcal{L}_u\{y'\} = sY - y(0^-)$, with \textbf{minus} in Laplace; unilateral Z shift is $z^{-1}Y + x[-1]$ with \textbf{plus}.
  \item \textbf{Applying final-value theorem without checking.} Requires $sX(s)$ (or $(1-z^{-1})X(z)$) to have all poles strictly in LHP / inside unit circle.
  \item \textbf{Mixing LHP with ``inside the unit circle.''} In CT stability is about $\operatorname{Re}\{s\}<0$. In DT it's about $|z|<1$. A DT pole at $z=-0.9$ is stable ($|-0.9|<1$) even though $z=-0.9$ is on the negative real axis.
  \item \textbf{Pole at $z=a$, not $z=1/a$.} For $1/(1-az^{-1})$, the pole is at $z=a$.
  \item \textbf{Working Z partial fractions in $z$, not $z^{-1}$.} All standard pairs are in $z^{-1}$ form. Use $z^{-1}$ as the variable.
  \item \textbf{Sampling at exactly $2\omega_M$.} Strictly insufficient (Example 22.3: $\sin(\omega_s t/2)$ samples to all zeros).
  \item \textbf{Confusing Hz and rad/s in sampling.} $f_s > 2f_M$ or $\omega_s > 2\omega_M$; $\omega = 2\pi f$.
  \item \textbf{Aliased frequency formula.} $f_\text{alias} = |f_s - f_0|$ only for $f_s/2 < f_0 < f_s$. More generally, aliasing folds components back into $[-f_s/2, f_s/2]$.
  \item \textbf{Sign error at the summer.} Negative feedback $\to 1+GH$; positive feedback $\to 1-GH$.
  \item \textbf{Open-loop vs.\ closed-loop poles.} Poles of $GH$ are open-loop; roots of $1+KGH=0$ are closed-loop.
  \item \textbf{Gain margin vs.\ phase margin confusion.} GM is read on magnitude plot at the frequency where phase $=-180^\circ$. PM is read on phase plot at the frequency where magnitude $=0$ dB.
  \item \textbf{Assuming feedback always stabilizes.} Too much gain can destabilize a stable plant. Feedback is a tool, not a panacea.
  \item \textbf{Forgetting $1/T$ scaling} in the sampled spectrum. Replicas are scaled by $1/T$; reconstruction filter compensates with gain $T$.
  \item \textbf{RHP zeros $\neq$ instability.} Only poles determine stability; RHP zeros give nonminimum-phase behavior but not instability.
\end{enumerate}

\vspace{1em}
\begin{center}
\emph{Prepared for CEC 315 Exam 3 --- Spring 2026.}
\end{center}

\end{document}
